% =====================================================================
\chapter{Soluciones de los ejercicios}
\label{ap:soluciones}
% =====================================================================

Cada solución se presenta en cinco partes: datos, planteamiento, desarrollo, resultado y comprobación. En los pasos intermedios se redondea a cuatro o cinco cifras significativas salvo que se indique otra cosa. Las longitudes van en milímetros salvo indicación contraria. Todos los valores numéricos se han comprobado con Python; los de los casos de estudio, con el cuaderno del módulo.

% ---------------------------------------------------------------------
\section{Tema 1}

\begin{solucion}{eje:1-repr}
\datos cuatro situaciones de uso.
\planteamiento aplicar el \cref{tab:repr} y la discusión de la sección~1.5.
\pasos
\begin{enumerate}[label=\arabic*.]
  \item (a) Abrir agujeros es un cambio de topología: hace falta una representación volumétrica (densidades por elemento, vóxeles, SDF o campo neuronal). Una malla con conectividad fija no puede.
  \item (b) Un motor dibuja triángulos: malla con UV, normales y niveles de detalle (\cref{tema:5}).
  \item (c) Un LiDAR mide puntos en la superficie, sin conectividad: nube de puntos (con normales estimadas si hacen falta).
  \item (d) PyTorch3D rasteriza mallas de forma derivable y ofrece pérdidas sobre ellas: malla deformable que parte de una esfera, optimizando las posiciones de los vértices (topología de esfera fija).
\end{enumerate}
\resultado \textbf{(a) volumétrica; (b) malla; (c) nube de puntos; (d) malla con conectividad fija}.
\comprobacion en (d) la forma final tendrá necesariamente $\chi=2$; si el objeto tuviera un asa, haría falta otra representación, coherente con (a).
\end{solucion}

\begin{solucion}{eje:1-tetra-cubo}
\datos tetraedro y cubo triangulado.
\planteamiento contar elementos y aplicar $\chi=\abs{V}-\abs{E}+\abs{F}$ y la \cref{prop:3f2e}.
\pasos
\begin{enumerate}[label=\arabic*.]
  \item Tetraedro con vértices $0,1,2,3$ y caras $(0,2,1)$, $(0,1,3)$, $(0,3,2)$, $(1,2,3)$: $\abs{V}=4$, $\abs{F}=4$ y aristas $\{0,1\}$, $\{0,2\}$, $\{0,3\}$, $\{1,2\}$, $\{1,3\}$, $\{2,3\}$, es decir, $\abs{E}=6$.
  \item $\chi=4-6+4=2$ y $2\abs{E}=12=3\abs{F}$.
  \item Cubo: 8 vértices; 6 caras cuadradas, cada una en 2 triángulos: $\abs{F}=12$. Aristas: 12 del cubo más 6 diagonales (una por cara cuadrada), $\abs{E}=18$.
  \item $\chi=8-18+12=2$ y $2\abs{E}=36=3\abs{F}$.
\end{enumerate}
\resultado \textbf{tetraedro: $(4,6,4)$; cubo triangulado: $(8,18,12)$; ambos con $\chi=2$}.
\comprobacion el \cref{cor:euler} predice $\abs{F}=2\abs{V}-4$: $2\cdot 4-4=4$ y $2\cdot 8-4=12$. Correcto.
\end{solucion}

\begin{solucion}{eje:1-euler}
\datos malla cerrada con $\chi=2$ y $\abs{F}=1280$; 12 bytes por vértice y 12 por cara.
\planteamiento \cref{prop:3f2e} y \cref{cor:euler}.
\pasos
\begin{enumerate}[label=\arabic*.]
  \item $\abs{E}=\tfrac32\abs{F}=1920$.
  \item $\abs{V}=(\abs{F}+4)/2=642$.
  \item Memoria: $642\cdot 12+1280\cdot 12=7704+15\,360=23\,064$ bytes.
\end{enumerate}
\resultado $\mathbf{\abs{E}=1920}$, $\mathbf{\abs{V}=642}$; \textbf{$23\,064$ bytes ($22{,}5$ KiB)}.
\comprobacion $\chi=642-1920+1280=2$. Una icosfera de nivel $k$ tiene $10\cdot 4^k+2$ vértices: con $k=3$, $642$.
\end{solucion}

\begin{solucion}{eje:1-sdf}
\datos $f_e(x)=\norm{x}-20$; caja de semilados $b=(10,10,5)$; puntos $p_1$, $p_2$, $p_3$.
\planteamiento evaluar las SDF; unión $=\min$, intersección $=\max$, diferencia $=\max(f_b,-f_e)$.
\pasos
\begin{enumerate}[label=\arabic*.]
  \item $p_1=(0,0,0)$: $f_e=-20$; $q=(-10,-10,-5)$, $\norm{\max(q,0)}=0$ y $\min(\max_iq_i,0)=-5$, luego $f_b=-5$. Unión $-20$; intersección $-5$; diferencia $\max(-5,20)=20$.
  \item $p_2=(15,0,0)$: $f_e=-5$; $q=(5,-10,-5)$, $\norm{\max(q,0)}=5$, $\min(5,0)=0$, luego $f_b=5$. Unión $-5$; intersección $5$; diferencia $\max(5,5)=5$.
  \item $p_3=(0,0,12)$: $f_e=-8$; $q=(-10,-10,7)$, $f_b=7$. Unión $-8$; intersección $7$; diferencia $\max(7,8)=8$.
\end{enumerate}
\resultado en el orden $(f_e,\ f_b,\ \text{unión},\ \text{intersección},\ \text{diferencia})$: \textbf{$p_1$: $(-20,-5,-20,-5,20)$}; \textbf{$p_2$: $(-5,5,-5,5,5)$}; \textbf{$p_3$: $(-8,7,-8,7,8)$}.
\comprobacion la caja está dentro de la esfera (su vértice más lejano está a $\sqrt{10^2+10^2+5^2}=15$), así que la diferencia caja menos esfera es vacía: su función es positiva en todos los puntos, como se obtiene.
\end{solucion}

\begin{solucion}{eje:1-mc}
\datos (a) $f=-0{,}3$ en $x=2$, $f=0{,}5$ en $x=3$. (b) $A=\{(0,0),(1,0)\}$, $B=\{(0,1),(2,0)\}$.
\planteamiento (a) interpolación lineal $t=f_0/(f_0-f_1)$; (b) \cref{def:chamfer}.
\pasos
\begin{enumerate}[label=\arabic*.]
  \item $t=-0{,}3/(-0{,}3-0{,}5)=0{,}375$; vértice en $x=2+0{,}375\cdot 1=2{,}375$ mm.
  \item De $A$ a $B$: para $(0,0)$, distancias al cuadrado $1$ y $4$, mínimo $1$; para $(1,0)$, $2$ y $1$, mínimo $1$. Media $1$.
  \item De $B$ a $A$: para $(0,1)$, $1$ y $2$, mínimo $1$; para $(2,0)$, $4$ y $1$, mínimo $1$. Media $1$.
  \item $\CD=1+1=2$.
\end{enumerate}
\resultado \textbf{(a) $x=2{,}375$ mm; (b) $\CD(A,B)=2$}.
\comprobacion la interpolación de $f$ en $2{,}375$ vale $-0{,}3+0{,}375\cdot 0{,}8=0$.
\end{solucion}

\begin{solucion}{eje:1-clamp}
\datos predicciones $(0{,}05;\ 0{,}30;\ -0{,}08;\ -0{,}40)$; verdades $(0{,}02;\ 0{,}12;\ 0{,}15;\ -0{,}02)$; $\delta=0{,}1$.
\planteamiento $\ell=\abs{\clip(\hat f,-\delta,\delta)-\clip(s,-\delta,\delta)}$ y media.
\pasos
\begin{enumerate}[label=\arabic*.]
  \item Punto 1: ningún recorte; $\abs{0{,}05-0{,}02}=0{,}03$.
  \item Punto 2: $\clip(0{,}30)=0{,}1$ y $\clip(0{,}12)=0{,}1$; pérdida $0$.
  \item Punto 3: $\clip(-0{,}08)=-0{,}08$ y $\clip(0{,}15)=0{,}1$; pérdida $0{,}18$.
  \item Punto 4: $\clip(-0{,}40)=-0{,}1$ y $-0{,}02$; pérdida $0{,}08$.
  \item Media: $(0{,}03+0+0{,}18+0{,}08)/4=0{,}0725$.
\end{enumerate}
\resultado \textbf{pérdida media $0{,}0725$}. El punto 2 no aporta: lejos de la superficie (a más de $\delta$), basta con acertar el signo y el valor exacto no se penaliza.
\comprobacion el punto 3 tiene el signo equivocado (dice dentro cuando está fuera) y es el que más pesa, como debe.
\end{solucion}

\begin{solucion}{eje:1-vol}
\datos tetraedro del \cref{ej:tetra}, $a=30$ mm; desplazamiento de $0{,}5$ mm de $v_1$ en $x$.
\planteamiento \cref{prop:volumen} con las caras que contienen $v_1$ en orden cíclico $(1,b,c)$.
\pasos
\begin{enumerate}[label=\arabic*.]
  \item Caras con el vértice 1: $(0,2,1)\to(1,0,2)$; $(0,1,3)\to(1,3,0)$; $(1,2,3)$.
  \item $\nabla_{v_1}\Vol=\tfrac16(v_0\times v_2+v_3\times v_0+v_2\times v_3)=\tfrac16\bigl(0+0+(0,a,0)\times(0,0,a)\bigr)=\tfrac16(a^2,0,0)$.
  \item Con $a=30$: $(150,0,0)$ mm².
  \item Base por altura: la cara opuesta a $v_1$ (plano $x=0$) tiene área $a^2/2=450$ mm²; $\partial\Vol/\partial x_1=450/3=150$ mm².
  \item $\Delta\Vol\approx 150\cdot 0{,}5=75$ mm³.
\end{enumerate}
\resultado $\mathbf{\nabla_{v_1}\Vol=(150,0,0)}$ \textbf{mm²}; \textbf{$\Delta\Vol=75$ mm³}.
\comprobacion el volumen es $x_1\cdot a\cdot a/6$, lineal en $x_1$: con $x_1=30{,}5$, $\Vol=30{,}5\cdot 900/6=4575$, exactamente $75$ más.
\end{solucion}

\begin{solucion}{eje:1-eikonal}
\datos $f=\norm{x-c}-r$; $g=\varphi\circ f$ con $\varphi'>0$ y $\varphi(0)=0$; $g(x)=\norm{x}^2-4$ en $(2,0,0)$.
\planteamiento derivar la norma; regla de la cadena (módulo~4, tema~2).
\pasos
\begin{enumerate}[label=\arabic*.]
  \item (a) Para $x\ne c$, $\nabla\norm{x-c}=(x-c)/\norm{x-c}$, de norma~1; la constante $r$ no influye.
  \item (b) $\varphi$ estrictamente creciente con $\varphi(0)=0$ cumple $\operatorname{signo}\varphi(t)=\operatorname{signo}t$, luego $\{g\le 0\}=\{f\le 0\}$ y $\{g=0\}=\{f=0\}$: el mismo sólido y la misma superficie.
  \item Por la regla de la cadena, $\nabla g=\varphi'(f)\nabla f$, y con $\norm{\nabla f}=1$, $\norm{\nabla g}=\varphi'(f)$.
  \item (c) $\norm{x}^2-4=(\norm{x}-2)(\norm{x}+2)=f(f+4)$ con $f=\norm{x}-2$, es decir, $\varphi(t)=t^2+4t$, creciente en el rango $t\ge-2$ de $f$. En la superficie, $\varphi'(0)=4$.
  \item Directamente, $\nabla g=2x=(4,0,0)$ en $(2,0,0)$, de norma~4; término eikonal $(4-1)^2=9$.
\end{enumerate}
\resultado \textbf{(c) $\norm{\nabla g}=4$ y término eikonal $9$}. La regularización empuja a la red hacia la reparametrización con $\varphi'=1$, la única que hace que los valores sean distancias.
\comprobacion para el punto $(3,0,0)$, a distancia~1 de la superficie, $g=5$ mientras que $f=1$: $g$ sobrestima la distancia, coherente con $\varphi'>1$.
\end{solucion}

\begin{solucion}{eje:1-traslacion}
\datos malla cerrada y orientada de forma coherente; traslación $t$.
\planteamiento desarrollar el producto vectorial como suma sobre aristas dirigidas; derivar respecto a $t$.
\pasos
\begin{enumerate}[label=\arabic*.]
  \item $(v_b-v_a)\times(v_c-v_a)=v_b\times v_c-v_b\times v_a-v_a\times v_c=v_a\times v_b+v_b\times v_c+v_c\times v_a$: es la suma de $u\times w$ sobre las tres aristas dirigidas $(u,w)$ de la cara.
  \item En una malla cerrada y coherente, cada arista $\{u,w\}$ aparece una vez como $(u,w)$ y otra como $(w,u)$, y $u\times w+w\times u=0$. Sumando sobre todas las caras, todo se cancela: (a) queda probado.
  \item Sea $S(t)=\tfrac16\sum_f\det(v_a+t,v_b+t,v_c+t)$. Por multilinealidad, $\nabla_tS=\tfrac16\sum_f\bigl[(v_b+t)\times(v_c+t)+(v_c+t)\times(v_a+t)+(v_a+t)\times(v_b+t)\bigr]$, que es la suma del paso~1 para la malla trasladada, nula por (a).
  \item $S$ tiene gradiente nulo en todo $t$, luego es constante: la fórmula no depende de la traslación.
  \item Trasladar todos los vértices es mover cada $v_i$ en $t$, así que $\nabla_tS\big|_{t=0}=\sum_i\nabla_{v_i}\Vol=0$.
\end{enumerate}
\resultado \textbf{$\sum_f(v_b-v_a)\times(v_c-v_a)=0$; el volumen es invariante por traslación y $\sum_i\nabla_{v_i}\Vol=0$}.
\comprobacion en el \cref{ej:tetra}, $\nabla_{v_3}\Vol=(0,0,150)$ y, por simetría, $\nabla_{v_1}=(150,0,0)$ y $\nabla_{v_2}=(0,150,0)$; para $v_0$, la cara opuesta tiene normal $(1,1,1)/\sqrt3$ y área $\tfrac{\sqrt3}{2}a^2$, de modo que $\nabla_{v_0}=-\tfrac13\cdot\tfrac{\sqrt3}{2}a^2\cdot(1,1,1)/\sqrt3=(-150,-150,-150)$. La suma es nula.
\end{solucion}

% ---------------------------------------------------------------------
\section{Tema 2}

\begin{solucion}{eje:2-vf}
\datos cuatro afirmaciones.
\planteamiento \cref{teo:convexo}, \cref{def:simp}, \cref{def:compliance} y \cref{def:filtros}.
\pasos
\begin{enumerate}[label=\arabic*.]
  \item (a) Falsa. Con $p=1$ la rigidez es proporcional al material y no hay penalización de los grises; el óptimo (convexo) suele tener grandes zonas intermedias.
  \item (b) Verdadera. Con $\Emin=0$, un elemento vacío no aporta rigidez y los nodos rodeados de vacío quedan sueltos: $K$ es singular.
  \item (c) Falsa. La flexibilidad $F^\T U$ es el trabajo de las cargas: pondera los desplazamientos en los puntos cargados; puede haber puntos sin carga con desplazamientos grandes.
  \item (d) Falsa. El filtro de sensibilidades es heurístico: el vector filtrado no es, en general, el gradiente de ninguna función; el filtro de densidades sí da el gradiente exacto del problema filtrado.
\end{enumerate}
\resultado \textbf{F, V, F, F}.
\comprobacion con el código del cuaderno, la viga MBB con $p=3$ tiene una medida de grises de $0{,}265$; con $p=1$ (mismo filtro) sube a $0{,}597$ y la flexibilidad baja a $165{,}3$ porque el material intermedio no se penaliza: el diseño gris no es fabricable.
\end{solucion}

\begin{solucion}{eje:2-gd}
\datos diseño generativo comercial.
\planteamiento qué omite el problema de flexibilidad y qué cambia el posprocesado (sección~2.7).
\pasos
\begin{enumerate}[label=\arabic*.]
  \item El diseño analizado durante la optimización es un campo de densidades en una malla gruesa; la pieza fabricada es una geometría suavizada y reconstruida, distinta.
  \item El objetivo es la rigidez global; no se controlan las tensiones locales (concentraciones en esquinas), el pandeo de barras esbeltas ni la fatiga.
  \item Los modelos de material y de fabricación (anisotropía de la impresión, tensiones residuales, soportes) no están en el problema.
\end{enumerate}
\resultado \textbf{hay que reanalizar porque la geometría cambia al posprocesar y porque el problema no incluye, por ejemplo, tensiones admisibles ni pandeo}.
\comprobacion la sección~2.7 indica que el flujo habitual termina con un análisis de elementos finitos más fino que incluye tensiones.
\end{solucion}

\begin{solucion}{eje:2-penal}
\datos $x_e=0{,}5$; $p=1,2,3$; $\Emin\approx 0$.
\planteamiento $E_e/(E_0x_e)=x_e^{p-1}$; resolver $x^3=0{,}5$.
\pasos
\begin{enumerate}[label=\arabic*.]
  \item $p=1$: $0{,}5^0=1$; $p=2$: $0{,}5$; $p=3$: $0{,}25$.
  \item Media rigidez con $p=3$: $x^3=0{,}5\Rightarrow x=0{,}5^{1/3}=0{,}7937$.
\end{enumerate}
\resultado \textbf{relaciones $1$; $0{,}5$; $0{,}25$; densidad necesaria $0{,}794$}.
\comprobacion $0{,}7937^3=0{,}5000$. Con $p=3$ hace falta el 79~\% del material para obtener el 50~\% de la rigidez: los grises son caros.
\end{solucion}

\begin{solucion}{eje:2-oc}
\datos $x=(0{,}5;0{,}5;0{,}5)$; $-\partial c/\partial x=(4;\ 1;\ 0{,}25)$; $\eta=1/2$; $m=0{,}2$; volumen objetivo $1{,}5$.
\planteamiento $x_e^{\text{nuevo}}=\clip\bigl(x_e\sqrt{-\partial_ec/\Lambda},\ [0{,}3;\ 0{,}7]\bigr)$; ajustar $\Lambda$.
\pasos
\begin{enumerate}[label=\arabic*.]
  \item Sin recortar, $x_e\sqrt{-\partial_ec/\Lambda}=(1;\ 0{,}5;\ 0{,}25)/\sqrt\Lambda$.
  \item El primero se recorta a $0{,}7$ si $1/\sqrt\Lambda\ge 0{,}7$, es decir, $\Lambda\le 2{,}04$, y el tercero a $0{,}3$ si $0{,}25/\sqrt\Lambda\le 0{,}3$, es decir, $\Lambda\ge 0{,}69$. Si ambos se recortan, la suma es $0{,}7+0{,}5/\sqrt\Lambda+0{,}3$, que debe valer $1{,}5$.
  \item $0{,}5/\sqrt\Lambda=0{,}5\Rightarrow\Lambda=1$, dentro de $[0{,}69;\ 2{,}04]$.
  \item $x^{\text{nuevo}}=(0{,}7;\ 0{,}5;\ 0{,}3)$.
\end{enumerate}
\resultado $\mathbf{\Lambda=1}$, $\mathbf{x^{\text{nuevo}}=(0{,}7;\ 0{,}5;\ 0{,}3)}$.
\comprobacion volumen $1{,}5$; el elemento más sensible gana material y el menos sensible lo pierde, ambos limitados por $m=0{,}2$. En el interior, $B_2=1/(\Lambda\cdot 1)=1$, la condición de optimalidad.
\end{solucion}

\begin{solucion}{eje:2-filtro}
\datos $\rmin=1{,}5$; sensibilidades: centro $-4$, laterales $-2$, diagonales $-1$; densidades uniformes.
\planteamiento \cref{def:filtros}; con $x$ uniforme, $x_i$ y $\max(\gamma,x_e)$ se cancelan y queda la media ponderada.
\pasos
\begin{enumerate}[label=\arabic*.]
  \item Pesos: centro $1{,}5-0=1{,}5$; laterales $1{,}5-1=0{,}5$; diagonales $1{,}5-\sqrt2=0{,}08579$. Los demás elementos están a distancia $\ge 2$ y pesan cero.
  \item Suma de pesos: $1{,}5+4\cdot 0{,}5+4\cdot 0{,}08579=3{,}8431$.
  \item Numerador: $1{,}5(-4)+4\cdot 0{,}5(-2)+4\cdot 0{,}08579(-1)=-6-4-0{,}3431=-10{,}3431$.
  \item Cociente: $-10{,}3431/3{,}8431=-2{,}6913$.
\end{enumerate}
\resultado \textbf{sensibilidad filtrada $-2{,}691$}.
\comprobacion es una media ponderada de valores entre $-4$ y $-1$, y queda entre ellos; el cuaderno comprueba que la suma de pesos es $3{,}8431$.
\end{solucion}

\begin{solucion}{eje:2-cadena}
\datos $H=\begin{psmallmatrix}1{,}5&0{,}5&0\\0{,}5&1{,}5&0{,}5\\0&0{,}5&1{,}5\end{psmallmatrix}$; $x=(1;0{,}5;0)$; $\partial c/\partial\xf=(-3;-2;-1)$.
\planteamiento $\xf_e=\sum_iH_{ei}x_i/\sum_iH_{ei}$ y $\partial c/\partial x_j=\sum_e(H_{ej}/\sum_iH_{ei})\,\partial c/\partial\xf_e$.
\pasos
\begin{enumerate}[label=\arabic*.]
  \item Sumas de filas: $(2;\ 2{,}5;\ 2)$.
  \item $\xf_1=(1{,}5+0{,}25)/2=0{,}875$; $\xf_2=(0{,}5+0{,}75+0)/2{,}5=0{,}5$; $\xf_3=(0{,}25+0)/2=0{,}125$.
  \item $\partial c/\partial x_1=\tfrac{1{,}5}{2}(-3)+\tfrac{0{,}5}{2{,}5}(-2)=-2{,}25-0{,}4=-2{,}65$.
  \item $\partial c/\partial x_2=\tfrac{0{,}5}{2}(-3)+\tfrac{1{,}5}{2{,}5}(-2)+\tfrac{0{,}5}{2}(-1)=-0{,}75-1{,}2-0{,}25=-2{,}2$.
  \item $\partial c/\partial x_3=\tfrac{0{,}5}{2{,}5}(-2)+\tfrac{1{,}5}{2}(-1)=-0{,}4-0{,}75=-1{,}15$.
\end{enumerate}
\resultado $\mathbf{\xf=(0{,}875;\ 0{,}5;\ 0{,}125)}$; $\mathbf{\partial c/\partial x=(-2{,}65;\ -2{,}2;\ -1{,}15)}$.
\comprobacion la suma se conserva, $-2{,}65-2{,}2-1{,}15=-6=-3-2-1$, porque cada fila de $H$ dividida por su suma suma~1. El volumen físico, $1{,}5$, coincide con el de diseño.
\end{solucion}

\begin{solucion}{eje:2-fisico}
\datos $\bar c=216{,}81$; $E_0=200$ GPa; $t=5$ mm; $F=2$ kN; semimodelo de $60\times 20$ mm al 50~\%; $7{,}85$ g/cm³.
\planteamiento $c=\bar cF^2/(E_0t)$, flecha $\delta=c/F$, masa $=\rho\cdot f_V\cdot$volumen.
\pasos
\begin{enumerate}[label=\arabic*.]
  \item $c=216{,}81\cdot(2000)^2/(2\cdot 10^{11}\cdot 0{,}005)=216{,}81\cdot 4\cdot 10^6/10^9=0{,}8672$ J.
  \item $\delta=0{,}8672/2000=4{,}336\cdot 10^{-4}$ m $=0{,}434$ mm.
  \item Volumen de material: $0{,}5\cdot 60\cdot 20\cdot 5=3000$ mm³ $=3$ cm³; masa $3\cdot 7{,}85=23{,}55$ g.
\end{enumerate}
\resultado \textbf{$c=0{,}867$ J; flecha $0{,}434$ mm; masa $23{,}6$ g}.
\comprobacion frente al aluminio del \cref{ej:mbb} ($0{,}310$ J con 1 kN): la carga doble multiplica por 4, $E_0t$ pasa de $7\cdot 10^8$ a $10^9$ N/m y $0{,}310\cdot 4\cdot 0{,}7=0{,}867$.
\end{solucion}

\begin{solucion}{eje:2-mecanismo}
\datos $g=L^\T U$; dos muelles en serie, $k=(512;\ 125)$ N/mm, $k_e=x_e^3k_0$, $k_0=1000$ N/mm, $x=(0{,}8;\ 0{,}5)$, $F=100$ N en el extremo; salida $g=u_1$.
\planteamiento \cref{teo:adjunto} con $\partial g/\partial U=L^\T$.
\pasos
\begin{enumerate}[label=\arabic*.]
  \item (a) Adjunto: $K\lambda=L$ ($K$ simétrica). Sensibilidad: $dg/dx_e=-\lambda^\T(\partial K/\partial x_e)U=-p\,x_e^{p-1}(E_0-\Emin)\,\lambda_e^\T k_0u_e$. Cuesta dos resoluciones: $U$ y $\lambda$, para cualquier número de variables.
  \item (b) $K=\begin{psmallmatrix}k_1+k_2&-k_2\\-k_2&k_2\end{psmallmatrix}$, $L=(1,0)^\T$. Se comprueba que $\lambda=(1/k_1,\ 1/k_1)$: primera fila $(k_1+k_2-k_2)/k_1=1$, segunda $(-k_2+k_2)/k_1=0$.
  \item $U=(F/k_1,\ F/k_1+F/k_2)$, así que $u_1=100/512=0{,}19531$ mm.
  \item Barra 1 (entre el apoyo y el nodo 1): $\partial K/\partial x_1=3x_1^2k_0\begin{psmallmatrix}1&0\\0&0\end{psmallmatrix}$, luego $dg/dx_1=-3x_1^2k_0\lambda_1u_1=-1920\cdot\tfrac{1}{512}\cdot 0{,}19531=-0{,}73242$ mm.
  \item Barra 2: $\partial K/\partial x_2\propto\begin{psmallmatrix}1&-1\\-1&1\end{psmallmatrix}$ y $dg/dx_2\propto-(\lambda_1-\lambda_2)(u_1-u_2)=0$.
\end{enumerate}
\resultado \textbf{$dg/dx_1=-0{,}7324$ mm y $dg/dx_2=0$}, con dos resoluciones.
\comprobacion directamente, $g=F/(x_1^3k_0)$, $dg/dx_1=-3F/(x_1^4k_0)=-300/(0{,}4096\cdot 1000)=-0{,}73242$ mm, y $g$ no depende de $x_2$.
\end{solucion}

\begin{solucion}{eje:2-kkt}
\datos problema de la \cref{def:compliance} con SIMP y cotas $0\le x_e\le 1$.
\planteamiento \cref{cor:autoadjunto}; condiciones de Karush--Kuhn--Tucker con multiplicadores $\Lambda\ge 0$ (volumen), $\mu_e^+\ge 0$ ($x_e\le 1$) y $\mu_e^-\ge 0$ ($x_e\ge 0$).
\pasos
\begin{enumerate}[label=\arabic*.]
  \item (a) Por el \cref{cor:autoadjunto}, $\partial c/\partial x_e=-p\,x_e^{p-1}(E_0-\Emin)u_e^\T k_0u_e\le 0$ porque $k_0$ es semidefinida positiva. Añadir material nunca empeora $c$.
  \item Si la restricción fuera $\sum x_e\le Nf_V$ y estuviera inactiva en un óptimo con algún $0<x_e<1$ y $\partial c/\partial x_e<0$, aumentar ligeramente $x_e$ seguiría siendo factible y reduciría $c$, contradicción. Luego la restricción como desigualdad es activa y equivale a la igualdad.
  \item (b) Lagrangiana: $\mathcal{L}=c+\Lambda(\sum_ex_e-Nf_V)+\sum_e\mu_e^+(x_e-1)-\sum_e\mu_e^-x_e$. Estacionariedad: $\partial c/\partial x_e+\Lambda+\mu_e^+-\mu_e^-=0$.
  \item Si $0<x_e<1$, holgura complementaria: $\mu_e^\pm=0$, luego $-\partial c/\partial x_e=\Lambda$, es decir, $B_e=1$.
  \item Si $x_e=1$: $\mu_e^-=0$ y $-\partial c/\partial x_e=\Lambda+\mu_e^+\ge\Lambda$, es decir, $B_e\ge 1$. Análogamente, si $x_e=0$, $B_e\le 1$.
\end{enumerate}
\resultado \textbf{la restricción de volumen es activa; $B_e=1$ en los elementos intermedios, $B_e\ge 1$ en los sólidos y $B_e\le 1$ en los vacíos}.
\comprobacion la actualización OC (\cref{def:oc}) no mueve un elemento con $B_e=1$, aumenta los de $B_e>1$ hasta la cota superior y disminuye los de $B_e<1$: sus puntos fijos son exactamente los puntos KKT.
\end{solucion}

% ---------------------------------------------------------------------
\section{Tema 3}

\begin{solucion}{eje:3-elegir}
\datos tres escenarios.
\planteamiento coste por gradiente: diferencias finitas $n+1$ simulaciones; adjunto o diferenciación automática, $O(1)$ si el código es derivable; sustituto, $M$ simulaciones iniciales.
\pasos
\begin{enumerate}[label=\arabic*.]
  \item (a) 3 parámetros, caja negra de 6 h: no se puede derivar el código; las diferencias finitas cuestan 4 simulaciones (24 h) por gradiente y son ruidosas. Un sustituto (proceso gaussiano o MLP) con unas decenas de simulaciones y verificación del óptimo.
  \item (b) 5000 densidades y código propio: el método adjunto (o reescribir el resolvedor en un marco derivable). Diferencias finitas costarían 5001 resoluciones por iteración; un sustituto en 5000 dimensiones necesitaría demasiadas muestras.
  \item (c) Choques con contactos: el simulador es discontinuo y la diferenciación automática da gradientes poco útiles; con 20 parámetros, un sustituto (u optimización sin derivadas) es más fiable.
\end{enumerate}
\resultado \textbf{(a) sustituto; (b) adjunto; (c) sustituto u optimización sin derivadas}.
\comprobacion coincide con el caso de la viga MBB (adjunto, 1200 variables) y con el de la aleta (2 variables: ambos métodos sirven).
\end{solucion}

\begin{solucion}{eje:3-pinn}
\datos aleta rectangular: $k\,t\,\theta''-2h\theta=0$, $\theta(0)=\theta_b$, $\theta'(L)=0$.
\planteamiento la pérdida PINN suma el residuo en puntos de colocación y las condiciones de contorno~\cite{pinn}.
\pasos
\begin{enumerate}[label=\arabic*.]
  \item Residuo: $\mathcal{L}_R=\frac1M\sum_{j=1}^M\bigl(k\,t\,\theta_\varphi''(x_j)-2h\,\theta_\varphi(x_j)\bigr)^2$ con $x_j\in(0,L)$; $\theta_\varphi''$ se obtiene por diferenciación automática respecto a la entrada.
  \item Contorno: $\mathcal{L}_B=\bigl(\theta_\varphi(0)-\theta_b\bigr)^2+\bigl(\theta_\varphi'(L)\bigr)^2$.
  \item Pérdida: $\mathcal{L}=\mathcal{L}_R+w_B\mathcal{L}_B$, con un peso $w_B$ que equilibra las escalas.
\end{enumerate}
\resultado \textbf{$\mathcal{L}=\frac1M\sum_j(kt\theta_\varphi''-2h\theta_\varphi)^2(x_j)+w_B\bigl[(\theta_\varphi(0)-\theta_b)^2+\theta_\varphi'(L)^2\bigr]$}. No hace falta malla porque los puntos de colocación pueden ser aleatorios y las derivadas se calculan sobre la red, no con diferencias entre nodos.
\comprobacion la solución exacta $\theta=\theta_b\cosh(m(L-x))/\cosh(mL)$ anula ambos términos.
\end{solucion}

\begin{solucion}{eje:3-aleta}
\datos $L=0{,}04$ m; $t=0{,}0015$ m; $h=60$; $k=200$; $\theta_b=50$ K.
\planteamiento $m=\sqrt{2h/(kt)}$, $q=\sqrt{2hkt}\,\theta_b\tanh(mL)$, $\eta=\tanh(mL)/(mL)$.
\pasos
\begin{enumerate}[label=\arabic*.]
  \item $m=\sqrt{120/(200\cdot 0{,}0015)}=\sqrt{400}=20$ m$^{-1}$; $mL=0{,}8$.
  \item $\tanh 0{,}8=0{,}66404$; $\sqrt{2hkt}=\sqrt{120\cdot 200\cdot 0{,}0015}=\sqrt{36}=6$ W/(m·K).
  \item $q=6\cdot 50\cdot 0{,}66404=199{,}21$ W/m.
  \item $\eta=0{,}66404/0{,}8=0{,}830$; Biot $=60\cdot 0{,}0015/400=2{,}25\cdot 10^{-4}$.
\end{enumerate}
\resultado \textbf{$m=20$ m$^{-1}$, $q=199{,}2$ W/m, $\eta=0{,}830$, $\mathrm{Bi}=2{,}3\cdot 10^{-4}$}.
\comprobacion el cuaderno da $q(0{,}04)=199{,}21$ W/m; está por debajo del óptimo ($220{,}4$ W/m en $L=58{,}6$ mm), porque esta aleta es demasiado corta y gruesa.
\end{solucion}

\begin{solucion}{eje:3-coste}
\datos $N=1200$ variables; 10 ms por resolución; 100 iteraciones.
\planteamiento diferencias hacia delante: $N+1$ resoluciones; adjunto: la del estado más una adjunta (cero si es autoadjunto).
\pasos
\begin{enumerate}[label=\arabic*.]
  \item Diferencias finitas: $1201\cdot 10$ ms $=12{,}01$ s por gradiente; 100 iteraciones, $1201$ s ($\approx 20$ min).
  \item Adjunto: 1 resolución (flexibilidad autoadjunta), 10 ms; en general 2, 20 ms. 100 iteraciones: entre 1 y 2 s.
  \item Cociente: $1201$ (o $600$ en el caso general).
\end{enumerate}
\resultado \textbf{12 s frente a 0,01--0,02 s por gradiente; 20 min frente a 1--2 s en total}.
\comprobacion el cuaderno resuelve la viga MBB en 106 iteraciones y unos 1,6 s, coherente con el adjunto.
\end{solucion}

\begin{solucion}{eje:3-implicita}
\datos $A(0)=\begin{psmallmatrix}2&-1\\-1&2\end{psmallmatrix}$; $b=(1,0)^\T$; $J=u_1+u_2$; $\partial A/\partial\theta=\begin{psmallmatrix}1&0\\0&0\end{psmallmatrix}$.
\planteamiento \cref{teo:implicita}: $\partial R/\partial\theta=(\partial A/\partial\theta)u$, adjunto $A^\T\lambda=(1,1)^\T$.
\pasos
\begin{enumerate}[label=\arabic*.]
  \item $\det A=3$; $u=A^{-1}b=\tfrac13\begin{psmallmatrix}2&1\\1&2\end{psmallmatrix}\begin{psmallmatrix}1\\0\end{psmallmatrix}=(2/3,\ 1/3)$; $J=1$.
  \item $A^\T\lambda=(1,1)^\T$: como $A(1,1)^\T=(1,1)^\T$, $\lambda=(1,1)$.
  \item $dJ/d\theta=-\lambda^\T(\partial A/\partial\theta)u=-\lambda_1u_1=-2/3$.
  \item Diferencia finita: con $\theta=0{,}01$, $\det=4{,}02-1=3{,}02$, $u=(2/3{,}02,\ 1/3{,}02)$, $J=3/3{,}02=0{,}99338$; pendiente $(0{,}99338-1)/0{,}01=-0{,}662$.
\end{enumerate}
\resultado $\mathbf{dJ/d\theta=-2/3\approx-0{,}667}$.
\comprobacion la diferencia finita hacia delante, $-0{,}662$, difiere en $O(\Delta\theta)$, como corresponde a un esquema de primer orden.
\end{solucion}

\begin{solucion}{eje:3-cuadratico}
\datos $(L,q)=(30;\ 165{,}377)$, $(60;\ 220{,}336)$, $(90;\ 198{,}201)$; óptimo verdadero $L=58{,}6$ mm, $q=220{,}42$ W/m.
\planteamiento con tres puntos, el ajuste por mínimos cuadrados interpola; para nodos equiespaciados ($h=30$), $\beta_2=(q_3-2q_2+q_1)/(2h^2)$ y el vértice está en $L^\star=L_2-h(q_3-q_1)/\bigl(2(q_3-2q_2+q_1)\bigr)$.
\pasos
\begin{enumerate}[label=\arabic*.]
  \item Segunda diferencia: $198{,}201-440{,}673+165{,}377=-77{,}095$; $\beta_2=-77{,}095/1800=-0{,}04283$ W/(m·mm²).
  \item $q_3-q_1=32{,}824$; $L^\star=60-30\cdot 32{,}824/(2\cdot(-77{,}095))=60+6{,}386=66{,}39$ mm.
  \item Máximo predicho: $q_2-(q_3-q_1)^2/\bigl(8(q_3-2q_2+q_1)\bigr)=220{,}336+1077{,}4/616{,}76=222{,}08$ W/m.
  \item Verificación con el simulador en $L^\star$: $218{,}13$ W/m.
\end{enumerate}
\resultado \textbf{el sustituto propone $L=66{,}4$ mm y predice $222{,}1$ W/m; el valor real es $218{,}1$ W/m, un 1,0~\% por debajo del óptimo verdadero y 4 W/m por debajo de lo prometido}.
\comprobacion la curva real es asimétrica (crece deprisa y decae despacio), y una parábola desplaza el vértice hacia la derecha. Se ilustra la \cref{prop:2eps}: el error del sustituto se traduce en pérdida de calidad, y la verificación lo revela.
\end{solucion}

\begin{solucion}{eje:3-recurrencia}
\datos $\theta_{k+1}=(1-a\Delta t)\theta_k$; $\theta_0=50$; $\Delta t=1$; $a=0{,}1$; $N=3$.
\planteamiento derivada de $\theta_3=\theta_0(1-a\Delta t)^3$; recurrencia adjunta con $\Phi(\theta,a)=(1-a\Delta t)\theta$, $\partial\Phi/\partial\theta=1-a\Delta t$, $\partial\Phi/\partial a=-\Delta t\,\theta$.
\pasos
\begin{enumerate}[label=\arabic*.]
  \item Estados: $\theta_1=45$; $\theta_2=40{,}5$; $\theta_3=36{,}45$ K.
  \item Cerrada: $d\theta_3/da=-3\Delta t\,\theta_0(1-a\Delta t)^2=-3\cdot 50\cdot 0{,}81=-121{,}5$ K·s.
  \item Adjunta: $\lambda_3=1$; $\lambda_2=0{,}9$; $\lambda_1=0{,}81$.
  \item $d\theta_3/da=\sum_{k=1}^3\lambda_k(-\Delta t\,\theta_{k-1})=-(1\cdot 40{,}5+0{,}9\cdot 45+0{,}81\cdot 50)=-(40{,}5+40{,}5+40{,}5)=-121{,}5$.
\end{enumerate}
\resultado $\mathbf{\theta_3=36{,}45}$ \textbf{K}; $\mathbf{d\theta_3/da=-121{,}5}$ \textbf{K·s}.
\comprobacion diferencia finita: con $a=0{,}101$, $\theta_3=50\cdot 0{,}899^3=36{,}329$; $(36{,}329-36{,}45)/0{,}001=-121{,}4$.
\end{solucion}

\begin{solucion}{eje:3-cota}
\datos $\sup\abs{\hat s-s}\le\varepsilon$; validación en tres puntos.
\planteamiento demostración de la \cref{prop:2eps}; contraejemplo con un sustituto que coincide en los puntos de validación.
\pasos
\begin{enumerate}[label=\arabic*.]
  \item (a) $s(\hat\theta)\le\hat s(\hat\theta)+\varepsilon\le\hat s(\theta^\star)+\varepsilon\le s(\theta^\star)+2\varepsilon$.
  \item (b) Sea $s(\theta)=\theta$ en $[0,1]$, con mínimo en $0$, y sea $\hat s(\theta)=\theta-10\,\phi(\theta)$, donde $\phi$ es un pulso continuo que vale~1 en $\theta=0{,}9$ y~0 fuera de $(0{,}85;\ 0{,}95)$.
  \item Validación en $\{0{,}2;\ 0{,}5;\ 0{,}8\}$: error nulo en los tres.
  \item El mínimo de $\hat s$ está en $0{,}9$ ($\hat s=-9{,}1$), donde $s=0{,}9$: el peor valor salvo un 10~\% del intervalo.
\end{enumerate}
\resultado \textbf{la cota exige error uniforme; un error de validación nulo en unos puntos no la garantiza}. El optimizador va precisamente a las zonas no validadas en que el sustituto es optimista; por eso se verifica el óptimo con el simulador y se añade a los datos.
\comprobacion en el ejemplo, $\varepsilon=\sup\abs{\hat s-s}=10$ y la cota da $s(\hat\theta)\le 20$, que se cumple ($0{,}9$): la proposición no se viola, solo es inútil si $\varepsilon$ es grande.
\end{solucion}

\begin{solucion}{eje:3-schmidt}
\datos $q(L)=\sqrt{2hkA/L}\,\theta_b\tanh(cL^{3/2})$, $c=\sqrt{2h/(kA)}$; (b) $h=40$, $k=170$, $A=5\cdot 10^{-5}$ m², $\theta_b=50$ K.
\planteamiento derivar $\ln q$ respecto a $L$ e igualar a cero; cambiar a la variable $\beta=cL^{3/2}$.
\pasos
\begin{enumerate}[label=\arabic*.]
  \item Con $t=A/L$, $m=\sqrt{2h/(kt)}=\sqrt{2hL/(kA)}=cL^{1/2}$ y $mL=cL^{3/2}=\beta$; además $\sqrt{2hkt}=\sqrt{2hkA/L}$. Esto da la expresión de $q$.
  \item $\ln q=\text{cte}-\tfrac12\ln L+\ln\tanh\beta$, y $d\beta/dL=\tfrac32\beta/L$.
  \item $\dfrac{d\ln q}{dL}=-\dfrac{1}{2L}+\dfrac{\operatorname{sech}^2\beta}{\tanh\beta}\cdot\dfrac{3\beta}{2L}=0\iff\tanh\beta=3\beta\operatorname{sech}^2\beta\iff\sinh\beta\cosh\beta=3\beta\iff\sinh 2\beta=6\beta$.
  \item (b) $c=\sqrt{80/(170\cdot 5\cdot 10^{-5})}=\sqrt{9411{,}8}=97{,}014$ m$^{-3/2}$.
  \item $L=(1{,}4192/97{,}014)^{2/3}=0{,}05982$ m; $t=A/L=0{,}836$ mm.
  \item $q=\sqrt{2\cdot 40\cdot 170\cdot 5\cdot 10^{-5}/0{,}05982}\cdot 50\cdot\tanh 1{,}4192=3{,}3717\cdot 50\cdot 0{,}8894=149{,}9$ W/m.
\end{enumerate}
\resultado \textbf{$\sinh 2\beta=6\beta$, $\beta=1{,}4192$; $L=59{,}8$ mm, $t=0{,}836$ mm, $q=149{,}9$ W/m}.
\comprobacion $\sinh(2{,}8384)=8{,}5154$ y $6\cdot 1{,}4192=8{,}5153$. Para $\beta$ pequeño, $\sinh 2\beta\approx 2\beta<6\beta$, y para $\beta$ grande, $\sinh 2\beta>6\beta$: hay una única raíz positiva, el máximo.
\end{solucion}

% ---------------------------------------------------------------------
\section{Tema 4}

\begin{solucion}{eje:4-gradcero}
\datos rasterizador clásico y pérdida por píxeles.
\planteamiento \cref{prop:grad-cero} y \cref{def:softras}.
\pasos
\begin{enumerate}[label=\arabic*.]
  \item Cada píxel vale 0 o 1 y es constante en un entorno de casi todas las configuraciones de vértices: el gradiente de la pérdida respecto a los vértices es nulo casi en todas partes.
  \item Por tanto, el descenso de gradiente no mueve la malla, aunque la silueta esté desplazada.
  \item SoftRas sustituye la indicadora por una sigmoide de la distancia al borde: los píxeles cercanos al contorno, dentro y fuera, dependen de forma suave de los vértices.
  \item Compromiso de $\sigma$: grande, gradientes desde lejos pero silueta borrosa y sesgada; pequeño, fidelidad pero gradientes solo en una franja estrecha.
\end{enumerate}
\resultado \textbf{el gradiente es nulo casi en todas partes; SoftRas lo hace informativo a cambio de un sesgo controlado por $\sigma$, que conviene reducir progresivamente}.
\comprobacion en el \cref{ej:siluetas}, el recocido de $\tau$ (el análogo de $\sigma$) eleva la $\IoU$ volumétrica de $0{,}744$ a $0{,}991$.
\end{solucion}

\begin{solucion}{eje:4-casco}
\datos taza con asa; 50 siluetas alrededor.
\planteamiento definición de casco visual~\cite{laurentini}.
\pasos
\begin{enumerate}[label=\arabic*.]
  \item Una silueta solo informa de qué rayos cortan el objeto. El casco visual es la intersección de los conos de visión de todas las siluetas: el mayor sólido compatible con ellas.
  \item El hueco interior de la taza está rodeado por la pared en todas las direcciones horizontales: ningún rayo de una vista lateral lo atraviesa sin cortar la pared. Desde arriba, el rayo que entra en el hueco corta el fondo, así que la silueta tampoco lo revela.
  \item El casco visual es, por tanto, una taza «maciza»; el asa sí se recupera, porque su agujero se ve de perfil desde algunas vistas.
  \item Información adicional: color o sombreado (consistencia fotométrica entre vistas, como en NeRF), mapas de profundidad o luz estructurada.
\end{enumerate}
\resultado \textbf{no; el interior es una concavidad que ninguna silueta ve, y hace falta información fotométrica o de profundidad}.
\comprobacion el agujero del asa es una característica de género (un asa, $g=1$) visible en las siluetas laterales, mientras que el hueco de la taza no cambia ninguna silueta: coherente con la definición.
\end{solucion}

\begin{solucion}{eje:4-softras}
\datos $\sigma=10^{-4}$; distancias $0{,}01$ (dentro), $0{,}01$ y $0{,}02$ (fuera); dos triángulos con $D=0{,}7311$ y $0{,}2689$.
\planteamiento $D=\sigm(\delta d^2/\sigma)$; $I=1-\prod(1-D_j)$.
\pasos
\begin{enumerate}[label=\arabic*.]
  \item Dentro, $d=0{,}01$: $d^2/\sigma=10^{-4}/10^{-4}=1$; $D=\sigm(1)=0{,}7311$.
  \item Fuera, $d=0{,}01$: $D=\sigm(-1)=0{,}2689$.
  \item Fuera, $d=0{,}02$: $d^2/\sigma=4$; $D=\sigm(-4)=0{,}01799$.
  \item Agregado: $1-(1-0{,}7311)(1-0{,}2689)=1-0{,}2689\cdot 0{,}7311=1-0{,}1966=0{,}8034$.
\end{enumerate}
\resultado \textbf{$D=0{,}731$; $0{,}269$; $0{,}0180$; silueta agregada $0{,}803$}.
\comprobacion $\sigm(1)+\sigm(-1)=1$; el agregado supera al mayor de los dos valores, como corresponde a la probabilidad de una unión.
\end{solucion}

\begin{solucion}{eje:4-volumen}
\datos $\delta_i=0{,}5$; $\sigma=(0{,}2;\ 1{,}0;\ 3{,}0)$; $c=(0{,}1;\ 0{,}8;\ 0{,}4)$; $t=(1{,}0;\ 1{,}5;\ 2{,}0)$.
\planteamiento cuadratura de NeRF: $\alpha_i=1-e^{-\sigma_i\delta_i}$, $T_i=\prod_{j<i}(1-\alpha_j)$, $w_i=T_i\alpha_i$.
\pasos
\begin{enumerate}[label=\arabic*.]
  \item $\sigma\delta=(0{,}1;\ 0{,}5;\ 1{,}5)$; $\alpha=(0{,}09516;\ 0{,}39347;\ 0{,}77687)$.
  \item $T=(1;\ e^{-0{,}1};\ e^{-0{,}6})=(1;\ 0{,}90484;\ 0{,}54881)$.
  \item $w=(0{,}09516;\ 0{,}35603;\ 0{,}42636)$.
  \item Color: $0{,}1\cdot 0{,}09516+0{,}8\cdot 0{,}35603+0{,}4\cdot 0{,}42636=0{,}00952+0{,}28482+0{,}17054=0{,}46488$.
  \item Profundidad: $\sum w_it_i=0{,}09516+0{,}53405+0{,}85271=1{,}4819$; normalizada por $\sum w_i$: $1{,}6887$.
\end{enumerate}
\resultado \textbf{pesos $(0{,}0952;\ 0{,}3560;\ 0{,}4264)$; color $0{,}465$; profundidad $1{,}482$ (o $1{,}689$ normalizada)}.
\comprobacion $\sum w_i=0{,}87754=1-e^{-2{,}1}$, como predice la \cref{prop:pesos}; el resto, $0{,}1225$, es la transmitancia hacia el fondo.
\end{solucion}

\begin{solucion}{eje:4-alfa}
\datos $o=(0{,}8;\ 0{,}6;\ 0{,}9)$; $P=(-0{,}5;\ -0{,}1;\ -2{,}0)$; colores rojo, verde y azul.
\planteamiento $\alpha_i=o_ie^{P_i}$, $T_i=\prod_{j<i}(1-\alpha_j)$, $C=\sum_iT_i\alpha_ic_i$.
\pasos
\begin{enumerate}[label=\arabic*.]
  \item $\alpha_1=0{,}8\cdot 0{,}60653=0{,}48522$; $\alpha_2=0{,}6\cdot 0{,}90484=0{,}54290$; $\alpha_3=0{,}9\cdot 0{,}13534=0{,}12180$.
  \item $T_1=1$; $T_2=0{,}51478$; $T_3=0{,}51478\cdot 0{,}45710=0{,}23530$.
  \item Contribuciones: $0{,}48522$ (R), $0{,}51478\cdot 0{,}54290=0{,}27947$ (G), $0{,}23530\cdot 0{,}12180=0{,}02866$ (B).
\end{enumerate}
\resultado $\mathbf{C=(0{,}485;\ 0{,}279;\ 0{,}029)}$.
\comprobacion opacidad acumulada $0{,}79336=1-T_4$, con $T_4=0{,}23530\cdot 0{,}87820=0{,}20664$. Aunque la tercera gaussiana es muy opaca, aporta poco porque está detrás y el píxel queda lejos de su centro.
\end{solucion}

\begin{solucion}{eje:4-pe}
\datos $p=0{,}25$; $L=2$; posición con $L=10$ y dirección con $L=4$.
\planteamiento $\gamma(p)=(\sin\pi p,\cos\pi p,\sin 2\pi p,\cos 2\pi p,\dots)$, aplicada a cada coordenada.
\pasos
\begin{enumerate}[label=\arabic*.]
  \item $\sin(\pi/4)=0{,}7071$; $\cos(\pi/4)=0{,}7071$; $\sin(\pi/2)=1$; $\cos(\pi/2)=0$.
  \item Posición: $3$ coordenadas $\times\,2$ funciones $\times\,10$ frecuencias $=60$.
  \item Dirección (vector unitario de 3 componentes): $3\cdot 2\cdot 4=24$.
\end{enumerate}
\resultado $\mathbf{\gamma(0{,}25)=(0{,}7071;\ 0{,}7071;\ 1;\ 0)}$; \textbf{60 y 24 componentes}.
\comprobacion cada par $(\sin,\cos)$ tiene norma~1: $0{,}7071^2+0{,}7071^2=1$ y $1^2+0^2=1$. La codificación es la misma idea que la posicional sinusoidal del Transformer (módulo~14, §4.2).
\end{solucion}

\begin{solucion}{eje:4-iou}
\datos $S=(0{,}9;\ 0{,}6;\ 0{,}2;\ 0)$; $T=(1;\ 1;\ 0;\ 0)$.
\planteamiento $I=\sum S_iT_i$, $U=\sum(S_i+T_i-S_iT_i)$; $\partial I/\partial S_i=T_i$, $\partial U/\partial S_i=1-T_i$.
\pasos
\begin{enumerate}[label=\arabic*.]
  \item $I=0{,}9+0{,}6=1{,}5$; $U=1+1+0{,}2+0=2{,}2$.
  \item $\IoU=1{,}5/2{,}2=0{,}6818$; pérdida $0{,}3182$.
  \item $\dfrac{\partial\IoU}{\partial S_3}=\dfrac{T_3U-I(1-T_3)}{U^2}=\dfrac{0-1{,}5}{4{,}84}=-0{,}3099$.
\end{enumerate}
\resultado \textbf{$\IoU=0{,}682$, pérdida $0{,}318$, $\partial\IoU/\partial S_3=-0{,}310$}. El signo negativo indica que el píxel 3, que está fuera del objetivo, debe apagarse.
\comprobacion para un píxel dentro del objetivo, $\partial\IoU/\partial S_1=(1\cdot 2{,}2-0)/4{,}84=0{,}4545>0$: hay que encenderlo más. Los signos son los esperados.
\end{solucion}

\begin{solucion}{eje:4-pesos-dem}
\datos cuadratura de NeRF con $\sigma_i\ge 0$.
\planteamiento suma telescópica; división de un intervalo en dos.
\pasos
\begin{enumerate}[label=\arabic*.]
  \item (a) Es la demostración de la \cref{prop:pesos}: $w_i=T_i-T_{i+1}$ y $\sum_i w_i=1-T_{N+1}\in[0,1]$.
  \item (b) Sea un intervalo de longitud $\delta$ y densidad $\sigma$ con transmitancia de entrada $T$. Sin dividir, su peso es $T(1-e^{-\sigma\delta})$.
  \item Dividido en dos mitades: $\alpha'=1-e^{-\sigma\delta/2}$; pesos $T\alpha'$ y $T(1-\alpha')\alpha'$; suma $T\bigl(1-(1-\alpha')^2\bigr)=T(1-e^{-\sigma\delta})$, el peso original.
  \item La transmitancia de salida es $T(1-\alpha')^2=Te^{-\sigma\delta}$, la misma que sin dividir, de modo que los pesos posteriores no cambian.
  \item (c) La fórmula del \emph{splatting} es $\sum_ic_i\alpha_i\prod_{j<i}(1-\alpha_j)=\sum_iT_i\alpha_ic_i$, la misma composición con $\alpha_i$ dado directamente por la gaussiana en lugar de por $1-e^{-\sigma_i\delta_i}$.
\end{enumerate}
\resultado \textbf{los pesos son no negativos y suman $1-T_{N+1}$; la cuadratura es coherente al refinar; NeRF y \emph{splatting} comparten la composición alfa}.
\comprobacion con los datos del \cref{eje:4-volumen}, la primera muestra ($\sigma\delta=0{,}1$) dividida en dos da $0{,}04877+0{,}95123\cdot 0{,}04877=0{,}09516$, igual que antes.
\end{solucion}

\begin{solucion}{eje:4-cov}
\datos $\Sigma=RSS^\T R^\T$, $R$ ortogonal, $S=\diag(s_1,s_2,s_3)$; $\Sigma'=JW\Sigma W^\T J^\T$.
\planteamiento escribir $\Sigma=MM^\T$ y usar $y^\T MM^\T y=\norm{M^\T y}^2$.
\pasos
\begin{enumerate}[label=\arabic*.]
  \item (a) Con $M=RS$, $\Sigma=MM^\T$, que es simétrica: $(MM^\T)^\T=MM^\T$.
  \item $y^\T\Sigma y=\norm{M^\T y}^2\ge 0$: semidefinida positiva.
  \item $y^\T\Sigma y=0\iff M^\T y=0$. Como $M^\T=SR^\T$ y $R^\T$ es invertible, $M^\T$ es invertible si y solo si $\det S=s_1s_2s_3\ne 0$. Así, $\Sigma$ es definida positiva si y solo si todas las $s_i$ son no nulas (sus autovalores son $s_i^2$ y sus autovectores, las columnas de $R$, por el teorema espectral del módulo~3).
  \item (b) $\Sigma'=(JWM)(JWM)^\T$, de la misma forma, luego semidefinida positiva.
  \item (c) Un paso de gradiente sobre las seis entradas libres puede dar una matriz con un autovalor negativo (por ejemplo, $\begin{psmallmatrix}1&2\\2&1\end{psmallmatrix}$ tiene autovalores $3$ y $-1$), que no es una covarianza; la gaussiana dejaría de estar definida.
\end{enumerate}
\resultado \textbf{la parametrización $RSS^\T R^\T$ garantiza covarianzas válidas por construcción; la directa no}.
\comprobacion con $R=I$ y $S=\diag(2,1,0{,}5)$, $\Sigma=\diag(4,1,0{,}25)$, definida positiva; con $s_3=0$ la gaussiana degenera en un disco plano.
\end{solucion}

% ---------------------------------------------------------------------
\section{Tema 5}

\begin{solucion}{eje:5-ms}
\datos $\nu=30,60,90,120,144$ Hz.
\planteamiento $1000/\nu$ ms (\cref{def:presupuesto}).
\pasos
\begin{enumerate}[label=\arabic*.]
  \item $33{,}33$; $16{,}67$; $11{,}11$; $8{,}33$; $6{,}94$ ms.
  \item Una media de 16,7 ms puede esconder fotogramas de 30 ms compensados por otros de 10 ms: cada fotograma que supera el presupuesto se pierde o se repite, y el jugador percibe un tirón.
\end{enumerate}
\resultado \textbf{$33{,}3$ / $16{,}7$ / $11{,}1$ / $8{,}3$ / $6{,}9$ ms; el presupuesto debe cumplirse en un percentil alto, no en la media}.
\comprobacion $16{,}67\cdot 60=1000$ ms.
\end{solucion}

\begin{solucion}{eje:5-quadric}
\datos planos $z=0$, $p_1=(0,0,1,0)$, y $x=0$, $p_2=(1,0,0,0)$; punto $(0{,}1;\ 0{,}2;\ 0{,}3)$.
\planteamiento $Q=p_1p_1^\T+p_2p_2^\T$ y $\Delta(v)=\bar v^\T Q\bar v$ (\cref{def:cuadrica}).
\pasos
\begin{enumerate}[label=\arabic*.]
  \item $Q=\diag(1,0,1,0)$.
  \item $\Delta=0{,}1^2+0{,}3^2=0{,}01+0{,}09=0{,}10$, la suma de las distancias al cuadrado a $x=0$ y $z=0$.
  \item El bloque $A=\diag(1,0,1)$ es singular porque las normales $(0,0,1)$ y $(1,0,0)$ no generan $\R^3$ (\cref{prop:qem}). Error nulo en toda la recta $x=z=0$ (el eje $y$, la arista entre los dos planos).
  \item El algoritmo no resuelve el sistema: evalúa el error en los extremos de la arista y en su punto medio y toma el menor~\cite[§4]{garland}.
\end{enumerate}
\resultado \textbf{$Q=\diag(1,0,1,0)$; $\Delta=0{,}10$; sin solución única (recta $x=z=0$); se recurre a los extremos o al punto medio}.
\comprobacion el vértice puede deslizarse por la arista del diedro sin cambiar la forma: es lo esperado en un borde recto.
\end{solucion}

\begin{solucion}{eje:5-tex}
\datos $2048\times 2048$ téxeles; RGBA8 (32 bits), BC5 (8 bits), BC1 (4 bits); mipmaps $\times 4/3$.
\planteamiento memoria $=2048^2\cdot\text{bpp}/8\cdot 4/3$ bytes.
\pasos
\begin{enumerate}[label=\arabic*.]
  \item $2048^2=4\,194\,304$ téxeles.
  \item RGBA8: $4\,194\,304\cdot 4\cdot 4/3=22\,369\,621$ B $=21{,}33$ MiB.
  \item BC5: un cuarto, $5{,}33$ MiB; BC1: un octavo, $2{,}67$ MiB.
\end{enumerate}
\resultado \textbf{$21{,}33$ / $5{,}33$ / $2{,}67$ MiB}. BC5 codifica dos canales independientes con 8 bytes cada uno: las componentes $x,y$ de la normal se guardan con precisión y $z$ se reconstruye con $z=\sqrt{1-x^2-y^2}$; BC1 obliga a que los tres canales compartan una única recta de color por bloque, lo que distorsiona las normales.
\comprobacion la cuenta coincide con la del cuaderno escalada por $1/4$ (allí $4096^2$): $85{,}33/4=21{,}33$.
\end{solucion}

\begin{solucion}{eje:5-bc1}
\datos $y=(10,40,70,95)$; $w=(0,\tfrac13,\tfrac23,1)$.
\planteamiento mínimos cuadrados con filas $(1-w_i,\ w_i)$: ecuaciones normales $M^\T Me=M^\T y$ (módulo~8, §2.2).
\pasos
\begin{enumerate}[label=\arabic*.]
  \item $M^\T M=\begin{psmallmatrix}\sum(1-w)^2&\sum w(1-w)\\ \sum w(1-w)&\sum w^2\end{psmallmatrix}=\begin{psmallmatrix}14/9&4/9\\4/9&14/9\end{psmallmatrix}$.
  \item $M^\T y=\bigl(10+\tfrac23 40+\tfrac13 70,\ \tfrac13 40+\tfrac23 70+95\bigr)=(60,\ 155)$.
  \item Multiplicando por 9: $14e_0+4e_1=540$ y $4e_0+14e_1=1395$. Restando las ecuaciones: $10(e_1-e_0)=855$, $e_1-e_0=85{,}5$; sumándolas: $18(e_0+e_1)=1935$, $e_0+e_1=107{,}5$.
  \item $e_0=11$, $e_1=96{,}5$. Reconstrucción: $(11;\ 39{,}5;\ 68;\ 96{,}5)$; residuos $(-1;\ 0{,}5;\ 2;\ -1{,}5)$.
  \item Error cuadrático: $1+0{,}25+4+2{,}25=7{,}5$.
\end{enumerate}
\resultado $\mathbf{e_0=11}$, $\mathbf{e_1=96{,}5}$; \textbf{error cuadrático $7{,}5$}.
\comprobacion los residuos son ortogonales a las columnas de $M$: $\sum(1-w_i)r_i=-1+\tfrac23 0{,}5+\tfrac13 2+0=0$ y $\sum w_ir_i=\tfrac13 0{,}5+\tfrac23 2-1{,}5=0$.
\end{solucion}

\begin{solucion}{eje:5-lod}
\datos $\varepsilon=2$ mm; $\phi_v=60^\circ$; $H=1080$ y $2160$; $\tau=1$ px.
\planteamiento \cref{prop:error-px}: $d^\star=\varepsilon H/(2\tau\tan 30^\circ)$.
\pasos
\begin{enumerate}[label=\arabic*.]
  \item $\tan 30^\circ=0{,}57735$.
  \item $H=1080$: $d^\star=0{,}002\cdot 1080/(2\cdot 0{,}57735)=1{,}871$ m.
  \item $H=2160$: el doble, $3{,}741$ m.
\end{enumerate}
\resultado \textbf{a partir de $1{,}87$ m en 1080p y de $3{,}74$ m en 2160p}.
\comprobacion a $1{,}871$ m la altura visible es $2\cdot 1{,}871\cdot 0{,}57735=2{,}160$ m en 1080 píxeles: 2 mm por píxel, justo el error.
\end{solucion}

\begin{solucion}{eje:5-mochila}
\datos grupo~1 (peso 1): $(2{,}0;0)$, $(1{,}2;1)$, $(0{,}5;3)$; grupo~2 (peso 2): $(1{,}5;0)$, $(0{,}8;0{,}5)$, $(0{,}4;2)$; $B=2{,}5$ ms.
\planteamiento enumerar las 9 combinaciones; $D(\lambda)=\sum_g\min_\ell(q_{g\ell}+\lambda t_{g\ell})-\lambda B$.
\pasos
\begin{enumerate}[label=\arabic*.]
  \item Pérdidas ponderadas del grupo~2: $0$; $1$; $4$.
  \item Factibles ($T\le 2{,}5$): $(0,2)$: 2,4 ms, $Q=4$; $(1,1)$: 2,0 ms, $Q=2$; $(1,2)$: 1,6, $Q=5$; $(2,0)$: 2,0, $Q=3$; $(2,1)$: 1,3, $Q=4$; $(2,2)$: 0,9, $Q=7$. Óptimo: $(1,1)$ con $Q^\star=2$.
  \item $\lambda=1{,}25$: en el grupo~1, los valores de $q+\lambda t$ son $2{,}5$, $2{,}5$ y $3{,}625$, con mínimo $2{,}5$; en el grupo~2, son $1{,}875$, $2{,}0$ y $4{,}5$, con mínimo $1{,}875$. Entonces $D=4{,}375-3{,}125=1{,}25$.
  \item $\lambda=10/7$: en el grupo~1, los valores de $q+\lambda t$ son $2{,}857$, $2{,}714$ y $3{,}714$, con mínimo $2{,}714$; en el grupo~2, son $2{,}143$, $2{,}143$ y $4{,}571$, con mínimo $2{,}143$ (empate). Entonces $D=4{,}857-3{,}571=1{,}286=9/7$.
\end{enumerate}
\resultado \textbf{óptimo $(\text{LOD1},\text{LOD1})$ con 2,0 ms y $Q^\star=2$; $D(1{,}25)=1{,}25$ y $D(10/7)=9/7\approx 1{,}29$}, ambas menores que $Q^\star$.
\comprobacion la \cref{prop:dual} se cumple: $1{,}25\le 2$ y $1{,}29\le 2$. El hueco de $5/7$ se debe a que la elección es discreta y el óptimo deja 0,5 ms sin usar.
\end{solucion}

\begin{solucion}{eje:5-shelf}
\datos atlas $1024\times 1024$; seis cartas.
\planteamiento ordenar por altura decreciente y llenar estantes de izquierda a derecha; la altura del estante es la de su primera carta.
\pasos
\begin{enumerate}[label=\arabic*.]
  \item Orden: $512\times 300$, $400\times 300$, $300\times 280$, $600\times 250$, $350\times 240$, $500\times 200$.
  \item Estante 1 (alto 300): $512+400=912$; la siguiente ($300$) no cabe ($1212>1024$).
  \item Estante 2 (alto 280): $300+600=900$; la siguiente ($350$) no cabe.
  \item Estante 3 (alto 240): $350+500=850$.
  \item Altura total: $300+280+240=820\le 1024$: caben.
  \item Área útil: $153\,600+120\,000+84\,000+150\,000+84\,000+100\,000=691\,600$; fracción $691\,600/1\,048\,576=0{,}660$.
\end{enumerate}
\resultado \textbf{caben en tres estantes; se aprovecha el 66~\% del atlas}.
\comprobacion queda una franja libre de $1024\times 204$ ($20$~\%) y huecos dentro de los estantes; una heurística más fina (\emph{skyline}) podría colocar la carta de $600\times 250$ en otro sitio.
\end{solucion}

\begin{solucion}{eje:5-qem}
\datos cuádricas de la \cref{def:cuadrica}.
\planteamiento demostración de la \cref{prop:qem}; linealidad de $\bar v^\T Q\bar v$ en $Q$.
\pasos
\begin{enumerate}[label=\arabic*.]
  \item (a)--(c) Como en la demostración de la \cref{prop:qem}.
  \item $\bar v^\T(Q_1+Q_2)\bar v=\bar v^\T Q_1\bar v+\bar v^\T Q_2\bar v=\sum_{p\in P_1}d_p(v)^2+\sum_{p\in P_2}d_p(v)^2$: la suma de distancias al cuadrado a todos los planos de ambos conjuntos. Los planos de las caras compartidas por los dos vértices se cuentan dos veces; el artículo lo acepta como aproximación~\cite[§5]{garland}.
  \item Como cada término es $\ge 0$, $\Delta_{Q_1+Q_2}(v)\ge\max\bigl(\Delta_{Q_1}(v),\Delta_{Q_2}(v)\bigr)$.
\end{enumerate}
\resultado \textbf{la cuádrica acumulada mide la distancia a la unión de los planos, y el error del colapso domina a cada error individual}. Por eso los costes crecen a lo largo de la simplificación y el montículo prioriza primero las zonas planas.
\comprobacion en el \cref{ej:esquina}, la esquina tiene error nulo en el origen; colapsar dos vértices de la misma cara plana también da error nulo, porque ambos comparten el plano.
\end{solucion}

\begin{solucion}{eje:5-dual}
\datos problema de la \cref{def:mckp}; \cref{ej:lod}.
\planteamiento demostración de la \cref{prop:dual}; mínimo de funciones afines.
\pasos
\begin{enumerate}[label=\arabic*.]
  \item La dualidad débil y el criterio de optimalidad se demuestran como en la \cref{prop:dual}.
  \item Para cada $e$ fija, $\lambda\mapsto Q(e)+\lambda(T(e)-B)$ es afín. $D$ es el mínimo de un número finito de funciones afines, luego cóncava (el mínimo puntual de funciones cóncavas es cóncavo).
  \item En el \cref{ej:lod}, la mejor cota ($2{,}30$) se alcanza en $\lambda=1{,}25$, cuya elección $(2,2,1,0)$ es factible pero deja holgura ($T=3{,}8<4$), así que $\lambda(T-B)=-0{,}25\ne 0$ y no se cumple la holgura complementaria.
  \item Ninguna elección consume exactamente 4,0 ms con menos pérdida: la discreción impide cerrar el hueco.
\end{enumerate}
\resultado \textbf{$D(\lambda)\le Q^\star$, $D$ es cóncava, y el hueco del ejemplo, $2{,}55-2{,}30=0{,}25$, se debe a la discreción}.
\comprobacion $D(1{,}25)=Q(e)+1{,}25(3{,}8-4)=2{,}55-0{,}25=2{,}30$, que coincide con el valor del cuaderno.
\end{solucion}

% ---------------------------------------------------------------------
\section{Tema 6}

\begin{solucion}{eje:6-checks}
\datos pieza para imprimir y recurso para motor.
\planteamiento sección~6.2.
\pasos
\begin{enumerate}[label=\arabic*.]
  \item Impresión: estanqueidad (cerrada, variedad, orientada); voladizos respecto a la dirección de impresión; espesor mínimo; volumen de impresión y unidades; cavidades cerradas.
  \item Motor: número de triángulos y materiales; UV presentes y sin solapes; normales coherentes; escala y pivote; niveles de detalle; nombres.
  \item Comunes: orientación coherente de las caras (normales hacia fuera), escala y unidades correctas, y ausencia de caras degeneradas.
\end{enumerate}
\resultado \textbf{comunes: orientación, escala y unidades, caras degeneradas}; la estanqueidad es imprescindible en impresión y conveniente, aunque no siempre exigida, en un motor.
\comprobacion la cara invertida del \cref{ej:validacion} afecta a ambos usos: volumen erróneo al imprimir y una cara negra o invisible en el motor.
\end{solucion}

\begin{solucion}{eje:6-sds}
\datos $\epsilon=(0{,}3;-0{,}2)$; $\hat\epsilon_c=(0{,}5;0{,}1)$; $\hat\epsilon_u=(0{,}2;-0{,}1)$; $s=7{,}5$; $w=1$; $x=\theta=(0{,}5;0{,}5)$; tasa $0{,}01$.
\planteamiento $\hat\epsilon=\hat\epsilon_u+s(\hat\epsilon_c-\hat\epsilon_u)$; gradiente $w(\hat\epsilon-\epsilon)\,\partial x/\partial\theta$ con $\partial x/\partial\theta=I$.
\pasos
\begin{enumerate}[label=\arabic*.]
  \item $\hat\epsilon_c-\hat\epsilon_u=(0{,}3;\ 0{,}2)$; $\hat\epsilon=(0{,}2+2{,}25;\ -0{,}1+1{,}5)=(2{,}45;\ 1{,}4)$.
  \item Gradiente: $(2{,}45-0{,}3;\ 1{,}4+0{,}2)=(2{,}15;\ 1{,}6)$.
  \item Paso: $\theta=(0{,}5-0{,}0215;\ 0{,}5-0{,}016)=(0{,}4785;\ 0{,}484)$.
\end{enumerate}
\resultado $\mathbf{\hat\epsilon=(2{,}45;\ 1{,}4)}$; \textbf{gradiente $(2{,}15;\ 1{,}6)$; $\theta=(0{,}4785;\ 0{,}484)$}.
\comprobacion sin guía ($s=1$), $\hat\epsilon=\hat\epsilon_c$ y el gradiente sería $(0{,}2;\ 0{,}3)$: la guía amplifica la dirección que distingue el texto de la imagen incondicional, como describe la tarea.
\end{solucion}

\begin{solucion}{eje:6-aristas}
\datos caras $(0,1,2)$, $(0,2,3)$, $(0,3,1)$, $(1,3,2)$, $(1,2,4)$.
\planteamiento contar caras por arista no dirigida (\cref{def:validacion}).
\pasos
\begin{enumerate}[label=\arabic*.]
  \item Las cuatro primeras caras forman un tetraedro estanco con aristas $\{0,1\}$, $\{0,2\}$, $\{0,3\}$, $\{1,2\}$, $\{1,3\}$ y $\{2,3\}$, cada una en dos caras.
  \item La quinta cara añade $\{1,2\}$ (que pasa a tener 3 caras), $\{2,4\}$ y $\{1,4\}$ (una cara cada una).
  \item Clasificación: 5 interiores, 1 no variedad ($\{1,2\}$) y 2 de borde ($\{2,4\}$, $\{1,4\}$). Total $\abs{E}=8$.
  \item $\chi=5-8+5=2$.
\end{enumerate}
\resultado \textbf{$\abs{E}=8$: 5 interiores, 1 no variedad, 2 de borde; $\chi=2$}. Defectos: una «aleta» triangular pegada a una arista del tetraedro, que hace la malla no variedad y abierta.
\comprobacion $\chi=2$ no garantiza nada: una malla defectuosa puede tener la característica de una esfera. Las comprobaciones locales de aristas son imprescindibles.
\end{solucion}

\begin{solucion}{eje:6-voladizo}
\datos áreas $(400,300,200,500,600)$ mm² con $n_z=(-1;\ -0{,}8;\ -0{,}6;\ 0;\ 0{,}5)$; la primera apoyada; umbral $-\cos 45^\circ=-0{,}7071$.
\planteamiento soporte si $n_z<-0{,}7071$ y la cara no está en la bandeja.
\pasos
\begin{enumerate}[label=\arabic*.]
  \item Cara de 400 mm²: apoyada, excluida.
  \item Cara de 300 mm²: $-0{,}8<-0{,}7071$, necesita soporte (forma $36{,}9^\circ$ con la horizontal).
  \item Cara de 200 mm²: $-0{,}6>-0{,}7071$, se autosoporta ($53{,}1^\circ$).
  \item Las demás miran a un lado o hacia arriba.
  \item Fracción: $300/2000=0{,}15$.
\end{enumerate}
\resultado \textbf{el 15~\% del área necesita soporte}.
\comprobacion el ángulo de una cara con la horizontal es el que forma su normal con la vertical descendente, $\arccos(-n_z)$: $\arccos 0{,}8=36{,}87^\circ<45^\circ$ (soporte) y $\arccos 0{,}6=53{,}13^\circ\ge 45^\circ$ (autosoportada).
\end{solucion}

\begin{solucion}{eje:6-mejorN}
\datos $p=0{,}2$; $1-\delta=0{,}9$; 30 s por generación y 2 s por validación.
\planteamiento \cref{prop:mejorN}; distribución geométrica para el número esperado.
\pasos
\begin{enumerate}[label=\arabic*.]
  \item $N\ge\ln 0{,}1/\ln 0{,}8=(-2{,}3026)/(-0{,}22314)=10{,}32$, luego $N=11$.
  \item $1-0{,}8^{11}=0{,}914$ (con 10, $0{,}893<0{,}9$).
  \item Tiempo del lote: $11\cdot(30+2)=352$ s, casi 6 min.
  \item De uno en uno, el número de intentos hasta el primer éxito es geométrico, de media $1/p=5$.
\end{enumerate}
\resultado \textbf{$N=11$; 352 s; 5 intentos de media}.
\comprobacion generar en serie y parar en el primer válido cuesta de media $5\cdot 32=160$ s, menos que el lote, pero con una cola larga: la probabilidad de necesitar más de 11 intentos es $0{,}8^{11}=0{,}086$.
\end{solucion}

\begin{solucion}{eje:6-mc}
\datos rejilla $128^3$ en $[-1{,}5;\ 1{,}5]^3$; lotes de 65\,536; 1 unidad $=50$ mm; pared de $0{,}8$ mm.
\planteamiento paso $=3/127$; número de nodos $128^3$.
\pasos
\begin{enumerate}[label=\arabic*.]
  \item Paso: $3/127=0{,}02362$ unidades.
  \item Evaluaciones: $128^3=2\,097\,152$; lotes: $2\,097\,152/65\,536=32$.
  \item En milímetros: $0{,}02362\cdot 50=1{,}18$ mm.
  \item Una pared de $0{,}8$ mm es más fina que el paso: puede caer entre dos nodos sin que ninguno tenga $f<0$ y desaparecer de la malla.
  \item Para que al menos un nodo quede dentro conviene un paso de la mitad del espesor, $0{,}4$ mm: $150/0{,}4=375$ intervalos, 376 nodos por eje ($5{,}3\cdot 10^{7}$ evaluaciones).
\end{enumerate}
\resultado \textbf{paso de $0{,}0236$ unidades ($1{,}18$ mm), $2{,}1\cdot 10^{6}$ evaluaciones en 32 lotes; la pared de $0{,}8$ mm no se resuelve}.
\comprobacion es el error (iv) de la \cref{obs:errores1}.
\end{solucion}

\begin{solucion}{eje:6-genero}
\datos malla cerrada, conexa, de género $g$; toro de $16\times 8$ vértices periódicos.
\planteamiento \cref{prop:3f2e} y $\chi=2-2g$.
\pasos
\begin{enumerate}[label=\arabic*.]
  \item (a) $\abs{E}=\tfrac32\abs{F}$, luego $\chi=\abs{V}-\tfrac32\abs{F}+\abs{F}=\abs{V}-\tfrac12\abs{F}=2-2g$.
  \item (b) $\abs{V}=128$; cuadriláteros $16\cdot 8=128$ (periodicidad: tantos como vértices), luego $\abs{F}=256$ y $\abs{E}=384$.
  \item $\chi=128-384+256=0$, de donde $g=1$.
  \item (c) Una arandela es un toro topológico: el validador comprueba que la malla es estanca, conexa y con $\abs{V}-\tfrac12\abs{F}=0$; si da 2, el agujero se ha cerrado; si da $-2$, hay un asa espuria.
\end{enumerate}
\resultado \textbf{(a) $\abs{V}-\tfrac12\abs{F}=2-2g$; (b) $(128,384,256)$, $\chi=0$, $g=1$; (c) exigir $g=1$}.
\comprobacion para el cubo triangulado ($\abs{V}=8$, $\abs{F}=12$), $8-6=2$ y $g=0$.
\end{solucion}

\begin{solucion}{eje:6-goodhart}
\datos generador de soportes a partir de texto; fabricación aditiva.
\planteamiento combinar la sección~6.3, la \cref{prop:mejorN} y la \cref{obs:errores6}.
\pasos
\begin{enumerate}[label=\arabic*.]
  \item \emph{Selección} (baratas, automáticas): estanqueidad y orientación (\cref{def:validacion}), escala y caja envolvente, volumen positivo, género esperado, fracción de área con voladizo y espesor mínimo por SDF.
  \item \emph{Comprobaciones independientes} (no usadas para seleccionar ni ajustar): autointersecciones, análisis de elementos finitos con las cargas del soporte (tensión máxima y flecha), simulación del laminado y revisión humana de una muestra aleatoria.
  \item \emph{Reparaciones admitidas}: soldar vértices duplicados, reorientar caras y cerrar agujeros menores que una tolerancia; nunca reparaciones que cambien cotas funcionales sin volver a medirlas.
  \item \emph{Política}: generar $N$ candidatos según la \cref{prop:mejorN} y elegir el mejor entre los que pasan la selección.
  \item \emph{Detección del engaño}: vigilar la tasa de fallos en las comprobaciones independientes entre los candidatos que pasan la selección. Si la tasa de aprobados en la selección sube mientras la de las independientes baja, el generador está explotando el validador (ley de Goodhart, módulo~21, §6.4); se amplía entonces la selección con la comprobación que falla y se renueva el conjunto independiente.
\end{enumerate}
\resultado \textbf{un validador en dos niveles (selección e independiente), reparaciones limitadas y seguimiento de la divergencia entre ambos niveles}.
\comprobacion el diseño responde a los errores (i), (iii), (iv) y (vi) de la \cref{obs:errores6}.
\end{solucion}
